% Módulo 10 — Apêndice C: glossário didático para leigos (R769)
\chapter{Glossário didático: cada princípio em linguagem simples}
\label{ap:glossario}
Este apêndice traduz cada termo técnico e cada lema para linguagem cotidiana, sem exigir matemática prévia. A sequência segue a ordem em que os conceitos aparecem no texto. Nota de leitura: cada ``sempre'' e ``nenhum'' abaixo traduz um lema verificado por máquina (Apêndice~A) e vale exatamente nos casos cobertos pelo lema correspondente.

\section{Termos}
\textit{Agente:} programa especializado que executa um tipo de tarefa. \textit{Roteador:} o mecanismo que escolhe qual agente chamar. \textit{Escore:} nota que mede quão adequado um candidato parece. \textit{Softmax:} regra que transforma notas em fatias de uma pizza que soma 100\%: nenhuma fatia negativa, nenhuma passa de 100\%, e o total fecha. \textit{Combinação convexa:} média ponderada cujos pesos somam 1 --- como misturar sucos em proporções que enchem exatamente o copo. \textit{Utilidade:} nota final combinada de um candidato. \textit{Argmax / top-1 / top-k:} escolher o melhor (ou os $k$ melhores) pela nota. \textit{Temperatura:} botão que afia ou suaviza as fatias: frio concentra no favorito, quente equaliza. \textit{Overflow:} estouro da calculadora com números gigantes; subtrair o máximo antes de exponenciar evita o estouro. \textit{Machine-checked:} verificado pelo computador até o nível dos axiomas, sem confiar em revisão humana. \textit{Sorry/admit:} atalhos proibidos que dizem ``confie em mim'' sem provar --- bloqueados aqui. \textit{Axioma:} regra básica aceita sem prova; este trabalho usa só três padrão da matemática formalizada: \textit{propext} (proposições iguais quando equivalentes), \textit{Classical.choice} (é possível escolher um elemento) e \textit{Quot.sound} (quocientes bem formados). \textit{Arrependimento (regret):} quanto a escolha perdeu em relação à melhor possível (0 = perfeito). \textit{Baseline:} alternativa simples de comparação. \textit{Ablação:} desligar uma peça para medir sua falta. \textit{Intervalo de confiança:} margem de incerteza da média medida. \textit{Semente:} número que torna sorteios repetíveis. \textit{Reprodutibilidade:} outra pessoa, mesmos passos, mesmos números.

\section{Os onze lemas, um por linha}
T1: as fatias sempre somam 100\%. T2/T6: a média ponderada nunca sai do intervalo entre o pior e o melhor ingrediente. T3: num grupo finito sempre há um melhor. T4: tirar o mesmo valor de todas as notas não muda as fatias. T5: cada fatia fica entre 0\% e 100\% havendo ao menos dois candidatos. T7: candidato único fica com 100\%. T8: nota maior, fatia maior. T9: os $k$ primeiros da fila ordenada superam todo o resto da fila. T10: esfriar aumenta a fatia do favorito. T11: expoentes deslocados nunca estouram.
